\begin{figure}
  \[
  \begin{array}{c}
    \infer{
      \tmcst s = \tmcst s \in \persort[\tmcst s'] \searrow R
    }{
      (\tmcst s : \tmcst s') \in \Ax
      &
      R \iff \persort
    }
    \qquad
    \infer{
      \dreflect {\tmcst s} e = \dreflect {\tmcst s} {e'} \in \persort \searrow R
    }{
      e = e' \in \perne
      &
      R \iff \perneu
    }
    \\[8pt]
    \infer{
      \dpi a B \rho = \dpi {a'} {B'} {\rho'} \in \persort[\tmcst s_3] \searrow R
    }{
      \begin{array}{c}
        (\tmcst s_1, \tmcst s_2, \tmcst s_3) \in \Ru
        \qquad
        a = a' \in \persort[\tmcst s_1] \searrow R_I
        \qquad
        \mhighlight {R_I \text{ is PER}}
        \\
        \forall \DD :: n = n' \in R_I.\ \evalexpfunc B {\denvext \rho n} = \evalexpfunc {B'} {\denvext {\rho'} {n'}} \in \persort[\tmcst s_2] \searrow R_O(\DD)
        \\
        \forall m, m' \in \Dom.\ m = m' \in R \iff (\forall \DD :: n = n' \in R_I.\ \evalappfunc m n = \evalappfunc {m'} {n'} \in R_O(\DD))
      \end{array}
    }
  \end{array}
  \]
  %
  \caption{Definition of the \acp{PER} for sorts}
  \label{fig:per-sorts}
\end{figure}
